% from ega1-6-fr.tex
\oldpage[I]{143}
une composante connexe $C$ de $X$ et distinct de $C$ ne peut être
irréductible, car cet ensemble n'étant pas connexe est réunion de deux
ensembles non vides disjoints à la fois ouverts et fermés dans $F$, donc
fermés dans $X$~; par suite c) entraîne b). On en conclut aussitôt que c)
entraîne d), deux composantes connexes distinctes étant sans point commun.

Nous n'avons pas utilisé jusqu'ici le fait que $X$ est localement noethérien.
Supposons maintenant cette hypothèse réalisée et montrons que d) entraîne
a)~: en vertu de (\hyperref[0.2.1.6-fr]{0, 2.1.6}), on peut se borner au cas
où l'espace $X$ est noethérien, donc n'a qu'un nombre fini de composantes
irréductibles. Comme celles-ci sont fermées et deux à deux disjointes, elles
sont ouvertes.

Enfin l'équivalence de d) et e) est valable sans supposer que l'espace
sous-jacent au préschéma $X$ soit localement noethérien. On peut en effet se
ramener au cas où $X=\operatorname{Spec}(A)$ est affine en vertu de
(\hyperref[0.2.1.6-fr]{0, 2.1.6})~; dire que $x$ n'est contenu que dans une
seule composante irréductible de $X$ signifie alors que $\mathfrak{j}_x$ ne
contient qu'un seul idéal minimal de $A$ (\hyperref[I.1.1.14-fr]{1.1.14}), ce
qui équivaut à dire que $\mathfrak{j}_x\mathscr{O}_x$ ne contient qu'un seul
idéal minimal de $\mathscr{O}_x$, d'où la conclusion.

\begin{corollary}[6.1.11]
\label{I.6.1.11-fr}
Soit $X$ un espace localement noethérien. Pour que $X$ soit irréductible, il
faut et il suffit que $X$ soit connexe et non vide, et que deux composantes
irréductibles distinctes de $X$ ne se rencontrent pas. Si $X$ est un
préschéma, cette dernière condition équivaut à ce que
$\operatorname{Spec}(\mathscr{O}_x)$ est irréductible pour tout $x\in X$.
\end{corollary}

La dernière partie a été vue dans (\hyperref[I.6.1.10-fr]{6.1.10})~; il n'y a
donc à prouver que la suffisance des conditions de la première assertion.
Mais d'après (\hyperref[I.6.1.10-fr]{6.1.10}), ces conditions entraînent que
les composantes irréductibles de $X$ sont ses composantes connexes, et comme
$X$ est connexe et non vide, il est irréductible.

\begin{corollary}[6.1.12]
\label{I.6.1.12-fr}
Soit $X$ un préschéma localement noethérien. Pour que $X$ soit intègre, il
faut et il suffit que $X$ soit connexe et que $\mathscr{O}_x$ soit intègre
pour tout $x\in X$.
\end{corollary}

\begin{proposition}[6.1.13]
\label{I.6.1.13-fr}
Soit $X$ un préschéma localement noethérien, et soit $x\in X$ un point tel que
le nilradical $\mathscr{N}_x$ de $\mathscr{O}_x$ soit premier (resp. que
$\mathscr{O}_x$ soit réduit, resp. intègre)~; alors il existe un voisinage
ouvert $U$ de $x$ qui est irréductible (resp. réduit, resp. intègre).
\end{proposition}

Il suffit de considérer les deux cas où $\mathscr{N}_x$ est premier et où
$\mathscr{N}_x=0$, la troisième hypothèse étant conjonction des deux
premières. Si $\mathscr{N}_x$ est premier, $x$ n'appartient qu'à une seule
composante irréductible $Y$ de $X$ (\hyperref[I.6.1.10-fr]{6.1.10})~; la
réunion des composantes irréductibles de $X$ ne contenant pas $x$ est fermée
(l'ensemble de ces composantes étant localement fini), et le complémentaire
$U$ de cette réunion est donc ouvert et contenu dans $Y$, donc irréductible
(\hyperref[0.2.1.6-fr]{0, 2.1.6}). Si $\mathscr{N}_x=0$, on a aussi
$\mathscr{N}_y=0$ pour tout $y$ dans un voisinage de $x$, car $\mathscr{N}$
est quasi-cohérent (\hyperref[I.5.1.1-fr]{5.1.1}), et par suite cohérent
puisque $X$ est localement noethérien, et la conclusion résulte de
(\hyperref[0.5.2.2-fr]{0, 5.2.2}).

\subsection{Préschémas artiniens.}
\label{subsection:I.6.2-fr}

\begin{definition}[6.2.1]
\label{I.6.2.1-fr}
On dit qu'un préschéma est \emph{artinien} s'il est affine et si son anneau
est artinien.
\end{definition}

